% ============================================================================
% ANEXO A — DERIVACIÓN DE LA SOLUCIÓN EXACTA DEL GBM VÍA LEMA DE ITÔ
% ============================================================================
\section{Derivación de la solución exacta del GBM}
\label{anexo:gbm}

Este anexo demuestra la fórmula~\eqref{eq:gbm-exacta} utilizada en §\ref{subsec:gbm}, partiendo de la ecuación diferencial estocástica~\eqref{eq:sde-gbm} y aplicando el lema de Itô. El nivel matemático presupuesto es el de un primer curso de cálculo estocástico; para una exposición completa se remite a Øksendal~\cite{oksendal2003}.

% ----------------------------------------------------------------------------
\subsection*{A.1\quad El lema de Itô}

Sea $\{W_t\}_{t \geq 0}$ un movimiento Browniano estándar sobre un espacio de probabilidad filtrado $(\Omega, \mathcal{F}, \{\mathcal{F}_t\}, \mathbb{P})$. Un proceso de Itô es cualquier proceso de la forma
\[
    dX_t \;=\; \mu_t\, dt + \sigma_t\, dW_t,
\]
donde $\mu_t$ y $\sigma_t$ son procesos adaptados suficientemente integrables.

\begin{teorema}[Lema de Itô, caso univariante]
\label{teorema:ito}
Sea $\{X_t\}$ un proceso de Itô con coeficientes $(\mu_t, \sigma_t)$ y sea $f: \mathbb{R}_{>0} \times \mathbb{R} \to \mathbb{R}$ una función de clase $\mathcal{C}^{1,2}$ (continua con derivadas parciales continuas hasta orden 1 en $t$ y orden 2 en $x$). Entonces el proceso $Y_t = f(t, X_t)$ es también un proceso de Itô y satisface
\begin{equation}
    dY_t \;=\; \frac{\partial f}{\partial t}(t, X_t)\,dt
            + \frac{\partial f}{\partial x}(t, X_t)\,dX_t
            + \frac{1}{2}\,\frac{\partial^2 f}{\partial x^2}(t, X_t)\,\sigma_t^2\,dt.
    \label{eq:lema-ito}
\end{equation}
\end{teorema}

La diferencia esencial respecto al cálculo ordinario es el término correctivo $\frac{1}{2}f_{xx}\sigma_t^2\,dt$, conocido como \emph{corrección de Itô}, que surge porque la variación cuadrática de $W_t$ es $d\langle W\rangle_t = dt$ (y no cero como ocurriría con una función de variación acotada ordinaria).

% ----------------------------------------------------------------------------
\subsection*{A.2\quad Aplicación a $\ln S_t$}

Partimos de la SDE del GBM:
\[
    dS_t \;=\; \mu S_t\,dt + \sigma S_t\,dW_t,
    \tag{\ref{eq:sde-gbm}}
\]
con $\mu, \sigma \in \mathbb{R}$, $\sigma > 0$. Definimos $f(t, x) = \ln x$ (independiente de $t$). Sus derivadas parciales son:
\[
    f_t = 0, \qquad f_x = \frac{1}{x}, \qquad f_{xx} = -\frac{1}{x^2}.
\]

Aplicando el Teorema~\ref{teorema:ito} con $X_t = S_t$, $\mu_t = \mu S_t$ y $\sigma_t = \sigma S_t$:
\begin{align*}
    d(\ln S_t)
    &= 0 \cdot dt + \frac{1}{S_t}\,dS_t + \frac{1}{2}\left(-\frac{1}{S_t^2}\right)(\sigma S_t)^2\,dt \\[4pt]
    &= \frac{1}{S_t}(\mu S_t\,dt + \sigma S_t\,dW_t) - \frac{\sigma^2}{2}\,dt \\[4pt]
    &= \left(\mu - \frac{\sigma^2}{2}\right)dt + \sigma\,dW_t.
\end{align*}

El proceso $\ln S_t$ satisface, por tanto, la SDE de un movimiento Browniano con deriva constante $\mu - \sigma^2/2$ y volatilidad constante $\sigma$. Esta es una ecuación diferencial estocástica \emph{lineal} cuyos coeficientes son deterministas.

% ----------------------------------------------------------------------------
\subsection*{A.3\quad Integración y solución exacta}

Integrando ambos miembros entre $t$ y $t + h$:
\[
    \ln S_{t+h} - \ln S_t
    \;=\; \left(\mu - \frac{\sigma^2}{2}\right)h + \sigma\,(W_{t+h} - W_t).
\]

El incremento del Browniano $W_{t+h} - W_t$ es independiente de $\mathcal{F}_t$ y tiene distribución $\mathcal{N}(0, h)$. Escribiéndolo como $\sqrt{h}\,Z$ con $Z \sim \mathcal{N}(0,1)$, y aplicando la función exponencial a ambos lados:

\begin{equation}
    S_{t+h} \;=\; S_t \exp\!\left(\left(\mu - \frac{\sigma^2}{2}\right)h + \sigma\sqrt{h}\,Z\right),
    \qquad Z \sim \mathcal{N}(0,1).
    \tag{\ref{eq:gbm-exacta}}
\end{equation}

Esta es la fórmula de simulación exacta empleada en la función \texttt{rollout} del programa (§\ref{subsec:rollout-cod}). Es \emph{exacta} en el sentido de que no introduce error de discretización: la distribución de $S_{t+h}$ condicionada a $S_t$ es exactamente log-normal con parámetros $\ln S_t + (\mu - \sigma^2/2)h$ y $\sigma^2 h$, independientemente del paso $h$ elegido.

% ----------------------------------------------------------------------------
\subsection*{A.4\quad Nota sobre la corrección $\sigma^2/2$}

El término $-\sigma^2/2$ en la exponente tiene una interpretación intuitiva. La media de $S_{t+h}$ condicionada a $S_t$ es
\[
    \mathbb{E}[S_{t+h} \mid S_t] \;=\; S_t\, e^{\mu h},
\]
que crece a la tasa $\mu$ (la deriva de la SDE original). Sin embargo, la media del \emph{logaritmo} del precio crece a la tasa $\mu - \sigma^2/2 < \mu$: la corrección de Jensen, $-\sigma^2/2$, recoge el coste de la convexidad de la función logaritmo frente a la distribución log-normal. Este es precisamente el término que la desviación~(D3) de §\ref{subsec:desviaciones} omite en la calibración práctica de $\hat\mu$.
